% Figure 1: method overview. Pure TikZ, no external assets.
% series -> context-only normalisation -> period-aligned frames -> visible past / masked future
%        -> video backbone -> predicted frames -> numeric read-out -> loss (train) / forecast (test)
\begin{figure}[t]
\centering
\begin{tikzpicture}[
  font=\scriptsize, x=1cm, y=1cm,
  box/.style={draw, rounded corners=1.5pt, align=center, inner sep=3.5pt},
  frm/.style={draw, minimum width=5.6mm, minimum height=5.6mm, inner sep=0pt, fill=black!6},
  fut/.style={frm, fill=fcred!10, draw=fcred, densely dashed},
  prd/.style={frm, fill=pixblue!10, draw=pixblue},
  flow/.style={-latex, semithick, draw=black!65},
  grad/.style={-latex, semithick, draw=fcred, densely dashed},
  outp/.style={-latex, semithick, draw=pixblue},
]

% ---------------- row A: series -> rendering -> clip
\node[box, align=center] (series) at (1.05,0) {observed series $x_{1:L}$\\[1pt]
  \tikz[baseline]{\draw[black!70, line width=0.5pt] plot[smooth, tension=0.7]
    coordinates {(0,0.05) (0.15,0.14) (0.3,0.02) (0.45,0.16) (0.6,0.06) (0.75,0.2) (0.9,0.1)};}};
\node[box, align=center, fill=black!2] (render) at (4.15,0)
  {context-only normalisation\\ period-aligned rendering\\ (one frame $=k$ periods)};
\draw[flow] (series) -- (render);

\begin{scope}[shift={(6.45,0)}]
  \foreach \i in {0,1,2,3} \node[frm] at (\i*0.63,0) {};
  \node[fut] at (4*0.63,0) {}; \node[fut] at (5*0.63,0) {};
  \draw[decorate] (0,0) -- (0,0);
  \node[anchor=north, font=\tiny] at (0.95,-0.36) {visible past};
  \node[anchor=north, font=\tiny, text=fcred] at (2.83,-0.36) {masked future};
\end{scope}
\draw[flow] (render.east) -- ++(0.25,0);

% ---------------- row B: backbone -> predicted -> read-out
\node[box, align=center, minimum height=8mm] (bb) at (1.5,-2.15) {video backbone\\ (VideoMAE-B)};
\begin{scope}[shift={(3.72,-2.15)}]
  \node[prd] at (0,0) {}; \node[prd] at (0.63,0) {};
  \node[anchor=north, font=\tiny, text=pixblue] at (0.32,-0.36) {predicted frames};
\end{scope}
\node[box, align=center, minimum height=8mm] (ro) at (6.55,-2.15)
  {differentiable\\ numeric read-out};
\node[box, align=center, draw=fcred, text=fcred] (loss) at (9.9,-1.72)
  {forecast-aligned loss\\ in value units};
\node[box, align=center, draw=pixblue, text=pixblue] (fc) at (9.9,-2.72)
  {forecast $\hat{x}_{L+1:L+H}$};

\draw[flow] (6.45+5*0.63+0.32,-0.30) -- ++(0,-0.55) -| (bb.north);
\draw[flow] (bb.east) -- (3.42,-2.15);
\draw[flow] (4.05,-2.15) -- (ro.west);
\draw[grad] (ro.east) -- ++(0.35,0) |- (loss.west);
\draw[outp] (ro.east) -- ++(0.35,0) |- (fc.west);
\draw[grad] (loss.north) -- ++(0,0.42) -| (bb.north east);
\node[font=\tiny, text=fcred, fill=white, inner sep=1pt] at (5.15,-1.30)
  {gradient, time-series continual pretraining only};

\node[anchor=west, font=\tiny, text=fcred]  at (0.05,-3.30) {$--$ training path};
\node[anchor=west, font=\tiny, text=pixblue] at (2.10,-3.30) {$-$ inference path};
\node[anchor=west, font=\tiny] at (4.20,-3.30)
  {no parameter update on the target dataset};

\end{tikzpicture}
\caption{The transfer framework. A series is normalised with the statistics of its visible context
only and rendered so that one frame holds $k$ periods, which puts value on the vertical axis and
phase on the horizontal one; the trailing frames are masked. A video backbone predicts those frames,
and a differentiable read-out decodes them back into values. During time-series continual
pretraining the loss is taken on the decoded values (Section~\ref{sec:objective}); at test time the
same read-out emits the forecast and no parameter is updated on the target dataset. The backbone
enters this pipeline with weights from video pretraining, which is the only stage that sees video.}
\label{fig:overview}
\end{figure}
